% Partie : méthodes
% Chapitre : démonstrations
% Section : schémas d'inférence

\documentclass[../../../include/open-logic-section]{subfiles}

\begin{document}
\olfileid[fr]{mth}{prf}{pat}

\olsection{Schémas d'inférence}

Une démonstration se compose d'inférences individuelles. Lorsque
nous faisons une inférence, nous le signalons généralement par un
mot comme «~donc~», «~ainsi~» ou «~par conséquent~». L'inférence
s'appuie souvent sur un ou deux faits dont nous disposons déjà :
il peut s'agir d'hypothèses ou de conclusions obtenues au cours de
la démonstration. Pour être clairs, nous pouvons numéroter ces
faits, puis indiquer, à chaque inférence, les énoncés utilisés.
Souvent, nous expliquons aussi \emph{pourquoi} la nouvelle
conclusion découle de ce qui précède. Certains schémas d'inférence
reviennent très souvent dans les démonstrations ; nous allons en
examiner quelques-uns. D'autres, comme la démonstration par
récurrence, sont plus élaborés et seront abordés plus loin.

Nous avons déjà rencontré un schéma d'inférence : déplier,
c'est-à-dire appliquer, une définition. Pour ce faire, nous
reformulons un énoncé qui contient le definiendum au moyen du
definiens. Supposons, par exemple, avoir déjà établi au cours
d'une démonstration que $D = E$ (a). Nous pouvons alors appliquer
la définition de $=$ pour les ensembles et conclure : «~Donc,
par définition et d'après (a), tout !!{element} de~$D$ est
!!a{element} de~$E$, et réciproquement.~»

Chose peut-être déroutante, il nous arrive de donner la
justification d'une inférence \emph{avant} de la faire plutôt
qu'au moment où nous la faisons. Supposons que nous n'ayons pas
encore démontré $D = E$, mais que ce soit notre objectif. En
appliquant la définition, nous pouvons reformuler cet objectif :
il faut démontrer que tout !!{element} de~$D$ est !!a{element}
de~$E$, et réciproquement. La démonstration aurait donc la forme
suivante : (a) démontrer que tout !!{element} de~$D$ est
!!a{element} de~$E$ ; (b) démontrer que tout !!{element} de~$E$
est !!a{element} de~$D$ ; (c) en déduire, de (a) et de (b), que
$D = E$ par définition. Cependant, nous ne l'écririons
généralement pas ainsi. Nous pourrions plutôt écrire :
\begin{quote}
Nous voulons montrer que $D = E$. Par définition de~$=$, cela
revient à montrer que tout !!{element} de~$D$ est !!a{element}
de~$E$, et réciproquement.

(a) \dots{} (démonstration que tout !!{element} de~$D$ est
!!a{element} de~$E$) \dots

(b) \dots{} (démonstration que tout !!{element} de~$E$ est
!!a{element} de~$D$) \dots
\end{quote}

\subsection{Utiliser une conjonction}

Le schéma d'inférence le plus simple consiste peut-être à tirer
d'une conjonction l'un de ses membres. Autrement dit, si nous avons
supposé ou déjà démontré $p$ et~$q$, nous pouvons conclure~$p$
(et aussi~$q$). Cette inférence est si élémentaire qu'elle reste
souvent implicite. Par exemple, après avoir déplié la définition
de $D = E$, nous avons établi que tout !!{element} de~$D$ est
!!a{element} de~$E$, et réciproquement. Nous pouvons en conclure
que tout !!{element} de~$E$ est !!a{element} de~$D$ : c'est le
sens de «~réciproquement~».

\subsection{Démontrer une conjonction}

Parfois, l'énoncé à démontrer est une conjonction : on vous demande
de «~démontrer $p$ et $q$~». Il suffit alors de faire deux choses :
démontrer $p$, puis démontrer $q$. Vous pouvez partager votre
démonstration en deux parties et les numéroter pour plus de clarté.
Dans vos premières notes, vous pourriez écrire «~(1) Démontrer
$p$~» en haut de la page, puis «~(2) Démontrer $q$~» au milieu.
Il arrive aussi que la conjonction ne figure pas explicitement
dans la question, mais apparaisse après analyse. Ainsi, si l'on
vous demande de démontrer $D = E$, le dépliage de la définition
de~$=$ montre qu'il faut établir que tout !!{element} de~$D$
est !!a{element} de~$E$ \emph{et} que tout !!{element} de~$E$
est !!a{element} de~$D$.

\subsection{Démontrer une disjonction}

Lorsque l'énoncé à démontrer est une disjonction, c'est-à-dire
qu'il a la forme «~$p$ ou $q$~», il suffit d'établir l'un des
deux membres. Toutefois, il n'arrive presque jamais que l'un
des membres découle immédiatement des hypothèses du théorème.
Le plus souvent, ces hypothèses sont elles-mêmes disjonctives,
ou bien on cherche à montrer que tous les objets d'un certain
type possèdent l'une de deux propriétés : certains possèdent
la première, les autres la seconde. C'est alors qu'une
démonstration par cas est utile (voir plus bas).

\subsection{Démonstration d'une implication}

De nombreux théorèmes ont la forme d'une implication : il faut
montrer que, si $p$ est vrai, alors $q$ l'est aussi. La mise en
place est simple : supposons l'antécédent de l'implication, ici
$p$, puis démontrons-en la conclusion~$q$. Si votre théorème dit
«~si $p$, alors $q$~», commencez donc par «~supposons $p$~» ;
à la fin, vous devez avoir démontré~$q$.

Une implication peut être formulée de plusieurs manières.
Outre «~si $p$, alors $q$~», on peut lire «~$p$ seulement si
$q$~», «~$q$ si $p$~» ou «~$q$, pourvu que $p$~». Toutes ces
formulations ont le même sens et demandent de supposer $p$
pour démontrer~$q$. Rappelons qu'une équivalence, «~$p$ si et
seulement si (ssi) $q$~», réunit deux implications : si $p$
alors $q$, et si $q$ alors~$p$. Il suffit donc de faire deux
démonstrations d'implication, une pour chaque sens. Parfois,
on peut cependant démontrer une équivalence en enchaînant
d'autres équivalences, de façon à partir de «~$p$~» pour
arriver à «~$q$~» ; il faut alors vérifier que \emph{chaque}
étape est bien une équivalence.

\subsection{Énoncés universels}

Utiliser un énoncé universel est simple : ce qui vaut pour
tous les objets vaut pour chacun d'eux en particulier.
Supposons, par exemple, que l'hypothèse d'une démonstration
soit $A \subseteq B$. En dépliant la définition de
$\subseteq$, cela signifie que, pour tout $x \in A$,
$x \in B$. Si nous savons aussi que $z \in A$, nous
pouvons donc conclure que $z \in B$.

Démontrer un énoncé universel peut sembler un peu plus délicat.
Ces énoncés se présentent souvent sous la forme «~si $x$
possède la propriété~$P$, alors il possède la propriété~$Q$~»
ou «~tous les $P$ sont des $Q$~». L'énoncé ne se présente
pas toujours exactement ainsi : il faut parfois s'exercer
à reconnaître ce que l'on nous demande de démontrer. Mais
il s'agit souvent d'établir que tous les objets qui possèdent
une propriété en possèdent une autre.

Pour démontrer un énoncé universel, on donne un nom, ou une
variable, à l'un des objets qui possèdent la première
propriété, puis on montre qu'il possède aussi la seconde.
Autrement dit, pour démontrer quelque chose concernant
\emph{tous} les $P$, on le démontre pour un $P$
\emph{arbitraire}. Le nom introduit désigne précisément
un tel objet. Nous utilisons habituellement une seule
lettre, choisie selon certaines conventions : $n$ pour
un nombre naturel, $!A$ pour une !!{formula}, $A$ pour
un ensemble, $f$ pour une fonction, etc.

Le point délicat est de conserver cette généralité pendant
toute la démonstration. On suppose qu'un objet arbitraire
(«~$x$~») possède la propriété~$P$, puis on établit qu'il
possède la propriété~$Q$, en n'utilisant que les définitions
et les hypothèses permises. Puisque nous n'avons rien
stipulé de particulier au sujet de $x$, hormis la propriété
$P$, nous pouvons affirmer que tout objet possédant $P$
possède aussi~$Q$. Ainsi, $x$ représente ici \emph{tous}
les objets qui ont la propriété~$P$.

\begin{prop}
  Pour tous ensembles $A$ et $B$, $A \subseteq A \cup B$.
\end{prop}

\begin{proof}
  Soient $A$ et $B$ des ensembles arbitraires. Nous voulons
  montrer que $A \subseteq A \cup B$. Par définition de
  $\subseteq$, il faut établir que, pour tout $x$, si
  $x \in A$, alors $x \in A \cup B$. Soit donc $x \in A$
  un !!{element} arbitraire de~$A$. Nous devons montrer
  que $x \in A \cup B$. Puisque $x \in A$, nous avons
  $x \in A$ ou $x \in B$. Donc
  $x \in \Setabs{z}{z \in A \lor z \in B}$.
  Par définition de $\cup$, cela signifie
  $x \in A \cup B$.
\end{proof}

\subsection{Démonstration par cas}

Supposons qu'une disjonction soit une hypothèse ou une
conclusion déjà établie : nous avons supposé ou démontré
que $p$ ou $q$ est vrai. Nous voulons démontrer~$r$.
Procédons en deux étapes : supposons d'abord $p$ et
démontrons~$r$ ; supposons ensuite $q$ et démontrons
à nouveau~$r$. Cela fonctionne parce que nous savons
ou supposons que l'une des deux possibilités se réalise.
Chaque étape établit qu'une possibilité suffit à rendre
$r$ vrai. Si les deux sont vraies, nous avons deux raisons,
et non une seule, de conclure~$r$ : il n'est pas nécessaire
de traiter séparément l'hypothèse conjointe $p$ et~$q$.
Pour annoncer cette méthode, nous disons que nous
«~distinguons des cas~».

Supposons, par exemple, savoir que $x \in B \cup C$.
L'ensemble $B \cup C$ est défini comme
$\Setabs{z}{z \in B \text{ ou } z \in C}$.
Autrement dit, par définition, $x \in B$ ou $x \in C$.
Pour en déduire $x \in A$, nous supposerions d'abord
$x \in B$ et démontrerions $x \in A$ sous cette
hypothèse ; nous supposerions ensuite $x \in C$ et
démontrerions de nouveau $x \in A$. Vous écririez
«~Distinguons deux cas~» sous l'hypothèse, puis
«~Cas (1) : $x \in B$~» et, plus bas,
«~Cas (2) : $x \in C$~». Il resterait à remplir
chacune des deux parties de la page.

La démonstration par cas est particulièrement utile
lorsque la conclusion est elle-même disjonctive.
En voici un exemple simple :

\begin{prop}
Supposons $B \subseteq D$ et $C \subseteq E$.
Alors $B \cup C \subseteq D \cup E$.
\end{prop}

\begin{proof}
  Supposons (a) $B \subseteq D$ et (b) $C \subseteq E$.
  Par définition, tout $x \in B$ appartient aussi à $D$
  (c), et tout $x \in C$ appartient aussi à $E$ (d).
  Pour montrer $B \cup C \subseteq D \cup E$, nous devons
  établir que, si $x \in B \cup C$, alors
  $x \in D \cup E$ (par définition de $\subseteq$).
  Or $x \in B \cup C$ si et seulement si $x \in B$
  ou $x \in C$ (par définition de~$\cup$).
  De même, $x \in D \cup E$ si et seulement si
  $x \in D$ ou $x \in E$. Il faut donc démontrer :
  pour tout $x$, si $x \in B$ ou $x \in C$,
  alors $x \in D$ ou $x \in E$.

  \begin{quote}
  Jusqu'ici, nous n'avons fait que déplier des définitions !{}
  Nous avons reformulé notre proposition sans
  $\subseteq$ ni $\cup$ : il reste à démontrer
  une implication universelle. Selon la méthode décrite
  plus haut, nous introduisons un objet $x$ arbitraire
  et supposons vraie la partie «~si~», puis nous
  démontrons la partie «~alors~». Autrement dit,
  nous supposons $x \in B$ ou $x \in C$ et montrons
  $x \in D$ ou $x \in E$.\footnote{
    Ce paragraphe explique seulement notre démarche :
    il ne fait pas partie de la démonstration et vous
    n'avez pas besoin de donner autant de détails dans
    vos propres démonstrations.}
  \end{quote}

  Supposons $x \in B$ ou $x \in C$. Nous devons
  montrer $x \in D$ ou $x \in E$. Distinguons
  deux cas.

  Cas 1 : $x \in B$. D'après (c), $x \in D$.
  Donc $x \in D$ ou $x \in E$. (Nous venons
  d'utiliser l'inférence étudiée dans la
  sous-section précédente !)

  Cas 2 : $x \in C$. D'après (d), $x \in E$.
  Donc $x \in D$ ou $x \in E$.
\end{proof}

\subsection{Démontrer un énoncé d'existence}

Lorsqu'on vous demande de démontrer un énoncé d'existence,
la question prend généralement la forme «~démontrer qu'il
existe un~$x$ tel que $\dots x \dots$~» : il faut établir
qu'un objet possède la propriété décrite par
«~$\dots x \dots$~». Il faut alors trouver un objet
approprié et montrer qu'il possède cette propriété.
Cela paraît simple, mais la démonstration peut être
délicate. Elle demande souvent de \emph{construire}
ou de \emph{définir} un objet, puis de vérifier qu'il
possède la propriété requise. Trouver le bon objet
peut être difficile ; vérifier sa propriété aussi ;
parfois, il n'est même pas facile de montrer que la
définition proposée désigne effectivement un objet !{}

En général, vous commencez par désigner l'objet :
«~soit $x$ \dots~», où les points de suspension
précisent l'objet envisagé. Au besoin, vous montrez
que cette description définit bien un objet existant.
Vous établissez ensuite que $x$ possède la
propriété~$Q$. Voici un exemple simple.

\begin{prop}
  Supposons $x \in B$. Alors il existe un~$A$
  tel que $A \subseteq B$ et $A \neq \emptyset$.
\end{prop}

\begin{proof}
  Supposons $x \in B$. Posons $A = \{x\}$.
  \begin{quote}
    Nous avons défini l'ensemble~$A$ en énumérant
    ses !!{element}s. Puisque $x$ est supposé être
    un objet et que nous pouvons former l'ensemble
    dont il est l'unique !!{element}, il n'est pas
    nécessaire ici de prouver séparément que notre
    définition donne bien un ensemble~$A$. Mais
    nous devons encore montrer que $A$ possède les
    propriétés exigées par la proposition : sans
    cela, la démonstration serait incomplète !{}
  \end{quote}
  Puisque $x \in A$, nous avons $A \neq \emptyset$.
  \begin{quote}
    Nous utilisons ici la définition $A = \{x\}$,
    ainsi que les faits évidents $x \in \{x\}$ et
    $x \notin \emptyset$.
  \end{quote}
  Puisque $x$ est l'unique !!{element} de~$\{x\}$
  et que $x \in B$, tout !!{element} de~$A$
  est aussi !!a{element} de~$B$. Par définition
  de~$\subseteq$, $A \subseteq B$.
\end{proof}

\subsection{Utiliser un énoncé d'existence}

Supposons que vous sachiez qu'un énoncé d'existence
est vrai, parce que vous l'avez démontré ou parce
qu'il est donné comme hypothèse : par exemple
«~pour un certain~$x$, $x \in A$~» ou
«~il existe un $x \in A$~». Pour l'utiliser dans
une démonstration, vous pouvez donner un nom à
l'un des objets dont il affirme l'existence.
Puisque $A$ contient au moins un objet, ce nom
peut en désigner un, même si vous ne savez pas
le choisir concrètement ni le décrire autrement
que par son appartenance à~$A$. Aux fins de la
démonstration, nommez-en donc un. Il est essentiel
de choisir un nom qui n'a pas encore servi et
qui n'apparaît pas dans vos hypothèses : sinon,
le raisonnement risque d'être erroné. Vous passez
ainsi de «~pour un certain $x$, $x \in A$~» à
«~soit $a \in A$~». Vous pouvez alors raisonner
sur~$a$, utiliser d'autres hypothèses, puis
aboutir à une conclusion~$p$. Si $a$ ne figure ni dans $p$
ni dans les hypothèses encore ouvertes autres que $a \in A$,
on peut conclure $p$ à partir de «~pour un certain $x$, $x \in A$~»
et des autres hypothèses utilisées.

\begin{prop}
Si $A \neq \emptyset$, alors $A \cup B \neq \emptyset$.
\end{prop}

\begin{proof}
  Supposons $A \neq \emptyset$. Il existe donc
  un $x$ tel que $x \in A$.
  \begin{quote}
    Nous avons d'abord reformulé l'hypothèse
    de la proposition. Cette hypothèse,
    $A \neq \emptyset$, cache un énoncé d'existence
    que l'on découvre en dépliant plusieurs
    définitions. Selon la définition de $=$,
    $A = \emptyset$ si et seulement si tout
    $x \in A$ appartient aussi à $\emptyset$
    et que tout $x \in \emptyset$ appartient
    aussi à $A$. En logique classique, en niant les deux
    membres de cette équivalence, on obtient :
    $A \neq \emptyset$ si et seulement si
    un certain $x \in A$ n'appartient pas
    à $\emptyset$, ou un certain
    $x \in \emptyset$ n'appartient pas à~$A$.
    Comme rien n'appartient à $\emptyset$,
    le second terme de la disjonction est
    toujours faux ; le premier se réduit
    à «~$x \in A$~». Ainsi,
    $A \neq \emptyset$ si et seulement s'il
    existe un $x$ tel que $x \in A$.
    Voilà un énoncé d'existence. Utilisons-le
    en donnant un nom à l'un des !!{element}s
    de~$A$ :
  \end{quote}
  Soit $a \in A$.
  \begin{quote}
    Nous avons nommé l'un des objets de~$A$.
    Nous allons raisonner sur~$a$, en veillant
    à ne rien supposer de plus que $a \in A$ :
  \end{quote}
  Puisque $a \in A$, nous avons
  $a \in A \cup B$ par définition de~$\cup$.
  Il existe donc un $x$ tel que $x \in A \cup B$,
  c'est-à-dire $A \cup B \neq \emptyset$.
  \begin{quote}
    Dans cette dernière étape, nous sommes
    passés de «~$a \in A \cup B$~» à
    «~il existe un $x$ tel que $x \in A \cup B$~».
    Ce dernier énoncé ne mentionne plus~$a$ :
    il découle donc de «~il existe un $x$ tel
    que $x \in A$~», sans dépendre du témoin
    choisi. Cela signifie que
    $A \cup B \neq \emptyset$.
  \end{quote}
\end{proof}

Il est peut-être préférable de distinguer les
variables liées, comme «~$x$~», des noms
hypothétiques de témoins, comme $a$, ainsi
que nous l'avons fait. En pratique, cependant,
nous ne le faisons pas toujours et employons
simplement $x$ :
\begin{quote}
Supposons $A \neq \emptyset$, c'est-à-dire
qu'il existe un $x \in A$. Par définition
de $\cup$, $x \in A \cup B$. Donc
$A \cup B \neq \emptyset$.
\end{quote}
Dans ce cas, il faut toutefois faire très
attention à employer des lettres distinctes
pour les témoins de différents énoncés
d'existence. Par exemple, ce qui suit
n'est \emph{pas} une démonstration correcte
de «~si $A \neq \emptyset$ et
$B \neq \emptyset$, alors
$A \cap B \neq \emptyset$~»,
énoncé qui est lui-même faux :
\begin{quote}
Supposons $A \neq \emptyset$ et
$B \neq \emptyset$. Il existe donc
un $x$ tel que $x \in A$, et il
existe aussi un $x$ tel que $x \in B$.
Puisque $x \in A$ et $x \in B$,
la définition de~$\cap$ donne
$x \in A \cap B$. Donc
$A \cap B \neq \emptyset$.
\end{quote}
Voyez-vous quelle étape est erronée
et pourquoi la conclusion ne vaut pas ?

\end{document}
